\section{Equivalence of Dissipative Channels}
\label{equivalence_of_maps}
\subsection{Continuous Channel}
Notice that the instantaneous kicks in the evolution generated by (\ref{Flqop}) take place at integer steps and we can allow for continuous dephasing between the kicks. For $t \in (m,m+1)$ we have 
\begin{align}
    \dot{\rho} = -i[H_\text{I}, \rho] + \mathcal{D}(\rho) \equiv \mathcal{L}_{\text{cont}}(\rho).
\end{align}
Then at $t=m$, we apply the kick
\begin{align}
\rho(m^+)=\UK\rho(m^-)\UK^\dagger, \quad \UK=e^{-i \HK}
\end{align}
Therefore the dynamical map between an integer period is given by 
\begin{align}
    \rho_{m+1}=\UK \exp^{\mathcal{L}_\text{cont}}(\rho_m)\UK^\dagger .
\end{align}
where $\mathcal{L}_{\text{cont}}$ is the continuous channel Liouvillian. We can therefore define the continuous channel dynamical map as the composition of continuous time dissipation followed by a kick 
\begin{align}
    \mathcal{M}_{\text{cont}}(\rho):=\mathcal{U}_\text{K} \circ e^{\mathcal{L}_{\text{cont}}}(\rho).
\end{align}
\subsection{Stroboscopic Channel}
For the stroboscopic channel we simply allow for the dissipation to accumulate over the integer time-period and do not resolve the micro-time within that period and so the system density matrix evolves via the following map
\begin{align}
    \mathcal{U}_{\text{sys}}(\rho)&= U_{\text{sys}}[\bm{h}] \rho U^\dagger_{\text{sys}}[\bm{h}] \nonumber\\
\end{align}
We then introduce a dissipation channel $\mathcal{E}_{\text{diss}}$ that acts once per period, and the complete dynamical map looks like
\begin{align}
    \mathcal{M}_{\text{strob}}(\rho):=\mathcal{E
    }_{\text{diss}} \circ \mathcal{U}_{\text{sys}}(\rho).
\end{align}
Such a stroboscopic channel closely resembles experimental setups, where drives are practically instantaneous and periodic and one only cares about integer stroboscopic times where Markovianity ensures that the environment is featureless on the scale of a period. So this can be thought of as a coarse-grained Floquet quantum channel.
\newline
In the weak dissipation limit $(\gamma \ll 1)$ both the continuous and the stroboscopic channel are equivalent up to $\mathcal{O}(\gamma)$ corrections.
\\[4pt]
Let's look at the continuous map
\begin{align}
    \mathcal{M}_{\text{cont}}= \mathcal{U}_{\text{K}} \circ e^{\mathcal{L}_{\text{cont}}}
\end{align}
where
\begin{align}
    \mathcal{L}_\text{cont} &\sim -i[\HI,\rho]+\gamma \mathcal{D}(\rho) \nonumber
\end{align}
Using BCH
\begin{align}
    e^{-i[\HI,\cdot]+\gamma \mathcal{D}(\cdot)}&= e^{-i[\HI,\cdot]}e^{\gamma \mathcal{D}(\cdot)} \exp \left(-\frac{\gamma}{2}[-i[\HI,\cdot],\mathcal{D}]+\cdots\right) \nonumber \\
    &\sim e^{-i[\HI,\cdot]}e^{\gamma \mathcal{D}(\cdot)} + \mathcal{O}(\gamma)
\end{align}
Plugging this into the map we have
\begin{align}
    \mathcal{M}_{\text{cont}}=\mathcal{U}_{\text{K}} \circ e^{-i[\HI,\cdot]} \circ e^{\gamma \mathcal{D}(\cdot)} + \mathcal{O}(\gamma).
\end{align}
Since the composition ${U}_{\text{K}} \circ e^{-i[\HI,\cdot]}$ is exactly $\mathcal{U}_{\text{sys}}=\mathcal{U}_{\text{K}}\circ \mathcal{U}_{I}$, the continuous map is 
\begin{align}
    \mathcal{M}_{\text{cont}}=\mathcal{U}_{\text{sys}} \circ e^{\gamma \mathcal{D}(\cdot)} + \mathcal{O}(\gamma) \nonumber\\
    \mathcal{M}_{\text{cont}}=\mathcal{M}_{\text{strob}} + \mathcal{O}(\gamma).
\end{align}
This is the exact same form as $\mathcal{M}_{\text{strob}}$, up to where you put the dissipation channel before or after the unitary channel. To first order in $\gamma$, the choice only differs by $\mathcal{O}(\gamma)$.